% ==============================================================================
% Doc. 262 — Acceptance without visualisation and inside-view of decompactified time
% (Merger of former documents 162 and 163)
% Johann Pascher, May 2026
% ==============================================================================

\chapter*{Doc.~262 -- Acceptance without visualisation and inside-view of decompactified time}
\addcontentsline{toc}{chapter}{Doc.~262 -- Acceptance and inside-view}

\section*{Document preface}

This document brings together, in two complementary parts, the
epistemic self-positioning of FFGFT regarding the visualisability of
$T^4$. \textbf{Part~A} clarifies the negative side: spatial
visualisability has, for over a century, not been an acceptance
criterion of physics, and FFGFT stands in this tradition rather than
outside it. \textbf{Part~B} adds the positive side: what a
deliberately delimited inside-view metaphor can deliver -- up to the
explicit framing of the brain-folds metaphor as an
\enquote{inside-view metaphor, not an ontological claim}. The two
parts belong together: without Part~A, Part~B would be misread as
esotericism; without Part~B, Part~A would remain purely defensive.

% ==============================================================================
\part*{Part A -- Acceptance without visualisation: On the epistemic status of $T^4$}
\addcontentsline{toc}{part}{Part A -- Acceptance without visualisation}
% ==============================================================================

\section*{Preliminary note to Part~A}

A theory that posits a four-dimensional identification torus $T^4$ as
the carrier of physics regularly encounters the objection: ``I cannot
visualise that.'' This part addresses the objection directly. It does
not claim to deliver a spatial visualisation of $T^4$ -- that is not
possible. It shows instead that \emph{visualisability} in the sense of
spatial intuition has, for over a century, \emph{not} been an
acceptance criterion of physics, and that FFGFT stands within this
methodological tradition, not outside it.

% ==============================================================================
\section{What is not visualisable is not unintelligible}
% ==============================================================================

Modern physics routinely works with mathematical constructs that are
not spatially visualisable, and accepts them because they work. A few
examples:

\begin{itemize}[leftmargin=2em,itemsep=4pt,topsep=4pt]
  \item \textbf{Imaginary numbers.} $\sqrt{-1}$ was rejected for
        centuries as ``impossible'' -- Descartes coined the
        derisive name ``imaginary'', which has stuck to this day.
        Today electrical engineering, signal processing and quantum
        mechanics use them as a matter of course. The decisive step
        towards acceptance was not spatial visualisation but Gauss's
        idea of representing them as a \emph{direction perpendicular}
        to the real axis -- a geometric extension that made the
        numbers tractable without ``picturing'' them in any everyday
        sense.

  \item \textbf{Complex-valued wave functions.} The wave function
        $\psi$ of quantum mechanics is complex-valued; its imaginary
        component is physically effective (interference, phase,
        entanglement). No one can spatially visualise $\psi$, and
        yet it is the central object of all of quantum theory.

  \item \textbf{Minkowski spacetime.} In 1908 Minkowski introduced
        time with a factor of $ict$ -- as an \emph{imaginary}
        coordinate -- to handle Lorentz transformations cleanly in
        geometric terms. Standard relativity therefore treats time
        formally not like a spatial dimension; all of contemporary
        high-energy and astrophysics builds on this construction.

  \item \textbf{Infinite-dimensional Hilbert spaces.} Any serious
        formulation of quantum mechanics lives in an
        infinite-dimensional, typically complex Hilbert space.
        Spatial visualisability? Out of the question. Acceptance?
        Complete.

  \item \textbf{Gauge groups, spinors, Lie algebras.} The Standard
        Model of particle physics formulates itself over
        $SU(3) \times SU(2) \times U(1)$. Spinors are not vectors in
        the usual sense -- they change sign under a $2\pi$ rotation,
        which seems intuitively absurd. Yet spinor mathematics is
        indispensable and its consequences are measured.

  \item \textbf{Complex manifolds, Calabi-Yau spaces.} String theory
        and parts of mathematical physics work with six- or
        ten-dimensional spaces whose geometry is not pictorial. They
        are accepted in their respective communities without anyone
        ``seeing'' them.

  \item \textbf{Negative numbers.} Rejected in medieval Europe
        because no one could picture ``$-3$ apples''. An apple is
        there or it is not; less than no apple is plainly nonsense.
        \emph{Three apples borrowed}, however, was obvious even
        then -- everyone understood what a debt is. The lesson of
        this story is not that negative numbers became more
        pictorial, but that the description itself had to change:
        not ``picture a $-3$-apple'' but \emph{describe a
        debt relation structurally}. As soon as the descriptive
        frame shifted, the matter was immediately accepted.
\end{itemize}

The list is not exhaustive; it is not even representative. But it
suffices to make the point: \emph{visualisability in the sense of
spatial intuition has, since the mathematisation of physics, not
been a necessary acceptance criterion.}

% ==============================================================================
\section{The actual acceptance criterion}
% ==============================================================================

If not spatial visualisability, then what? In the practice of physics
since roughly Gauss and Minkowski, two criteria are established:

\begin{enumerate}[leftmargin=2em,itemsep=4pt,topsep=4pt]
  \item \textbf{Internal consistency.} The mathematical structures
        must be contradiction-free. Theorems must follow from
        definitions and axioms; derived quantities must behave in
        well-defined ways.

  \item \textbf{Testable consequences.} The construction must
        deliver statements that can be compared with observation or
        measurement -- and stand or fall in that comparison.
\end{enumerate}

Both criteria are \emph{structural}, not visual. They concern the
mathematical and empirical behaviour of the theory, not its
appearance before the reader's inner eye. A theory that satisfies
both criteria is physically acceptable, regardless of whether its
subject can be drawn.

\textbf{This shift in the acceptance criterion is not mere
convention.} It is the methodological consequence of recognising
that our visual system is evolutionarily tuned to a three-dimensional
habitat and becomes systematically unreliable for structural claims
beyond that habitat. Visualisability is a \emph{psychological}
phenomenon, not an epistemic one.

% ==============================================================================
\section{$T^4$ as a structural outside perspective}
% ==============================================================================

What the constructs listed in §\,1 have in common is more than their
unvisualisability. They are \emph{structural outside perspectives} on
systems that, from any single inside view, would be only partially
visible.

\begin{itemize}[leftmargin=2em,itemsep=4pt,topsep=4pt]
  \item No observer sees \emph{Lorentz geometry} -- each sees only
        their own slice through it. But the Minkowski description
        gives us the geometry \emph{as a whole}, within which the
        various observer perspectives appear as particular slices.

  \item No experimenter sees the \emph{wave function} -- each
        measurement yields only a single outcome. But the
        Hilbert-space description describes the wave function \emph{as
        an object}, whose properties (norm, phase, superposition)
        stand the test of many measurements statistically.

  \item No one sees a \emph{gauge group}. But the Standard-Model
        description captures the internal structure of the
        interactions \emph{as group theory}, whose consequences
        (mass ratios, mixing angles, conservation laws) are
        measurable.
\end{itemize}

In each of these cases the lesson is the same: the mathematical
\emph{outside view} is richer than any possible inside experience.
It makes visible what cannot be visible from within the system --
and that is precisely its value.

\textbf{$T^4$ stands in this row.} It is not a spatial picture but a
structural outside perspective on the carrier of physics in which
quantum mechanics and relativity do not contradict each other. From
any inside perspective -- that is, from any concrete reference
frame -- QM and GR are incompatible; from the outside perspective
$T^4$ they are different slices of the same closed structure. This
is methodologically the same programme that Minkowski carried out
for time, Hilbert for quantum mechanics, and the gauge theorists for
internal symmetry: \emph{not to produce an intuition, but to find an
outside language}.

% ==============================================================================
\section{What this is not}
% ==============================================================================

Three demarcations, so that the claim is not overstretched:

\textbf{First: no privileged outside position.} ``Outside view'' here
means a different \emph{descriptive perspective}, not a standpoint
``outside physics''. We are never really outside; we change the
descriptive frame. The mathematical outside view is not \emph{the}
outside view (no such thing exists) but \emph{an} outside view -- one
that makes visible structures that would remain hidden from the
inside.

\textbf{Second: no ontological claim.} FFGFT does not claim ``the
world is really a $T^4$''. It claims: \emph{if we want to describe
the physical system from an outside perspective in which the
contradictory inside views of QM and GR can be unified, then $T^4$ is
the minimal structure that achieves this.} This is a substantially
weaker and at the same time more defensible claim than ``this is how
the world is''.

\textbf{Third: no demand for visualisation.} The aim of this document
is not to produce an intuition of $T^4$. That would be an empty
promise. The aim is to separate the acceptance question from the
visualisation question and to refer it back to the two criteria on
which modern physics actually operates: consistency and testable
consequences.

% ==============================================================================
\section{Consequence for the discussion of FFGFT}
% ==============================================================================

Anyone who rejects FFGFT on the grounds that $T^4$ is not visualisable
would, to be consistent, also have to reject:
\begin{itemize}[leftmargin=2em,itemsep=2pt,topsep=2pt]
  \item the complex-valued wave function (not visualisable),
  \item the Lorentz transformation with imaginary time (not
        visualisable),
  \item the infinite-dimensional Hilbert spaces (not visualisable),
  \item the spinors with $4\pi$ periodicity (not visualisable),
  \item the gauge groups of the Standard Model (not visualisable).
\end{itemize}

Since no one seriously demands this -- and modern physics would
collapse without these constructs -- visualisability is not the
decisive criterion in the case of $T^4$ either. The decisive criterion
is whether FFGFT meets the two standard requirements:

\begin{enumerate}[leftmargin=2em,itemsep=2pt,topsep=2pt]
  \item \textbf{Internal consistency.} Developed in Doc.~001a
        (foundations), Doc.~181 (4D torus geometry), Doc.~193
        (Non-Closure Theorem), Doc.~230 (Hilbert-space bridge).
  \item \textbf{Testable consequences.} Concrete predictions for
        lepton masses, $\alpha$, CHSH deviation, cosmological
        scales; some measured (Bell violation 96.9\,\% of the
        Tsirelson bound, May 2026, Doc.~147 §\,8), others currently
        below the measurement noise floor.
\end{enumerate}

These two questions are the scientifically legitimate ones. The
question of visualisability is not.

% ==============================================================================
\section{Summary}
% ==============================================================================

Visualisability in the sense of spatial intuition has, for over a
century, not been an acceptance criterion in physics. Internal
consistency and testable consequences have taken its place. The
mathematical constructs of modern physics -- imaginary numbers,
complex wave functions, Minkowski geometry, Hilbert spaces, gauge
groups, spinors -- are none of them spatially visualisable; they are
nonetheless accepted because they supply structural outside
perspectives on systems that, from any inside view, would remain
incomplete.

$T^4$ stands in this row. It is not visualisable -- and it need not
be. It is precisely defined, internally consistent, and its
consequences are derivable; some are already measured, others lie
below the current measurement noise. That is the methodological
situation. Anyone who rejects FFGFT for lack of visualisability
rejects not FFGFT but the methodological programme of modern physics
as a whole.

Imagination need not be visual to be effective.



% ==============================================================================
\part*{Part B -- Time as mode, mass as deformation, fractality as inside-view of decompactified time}
\addcontentsline{toc}{part}{Part B -- Time, mass, fractality}
% ==============================================================================

\section*{Preliminary note to Part~B}

Part~A established that $T^4$ is not spatially visualisable and need
not be. That is the negative side: \emph{what visualisation cannot
deliver}.

This part makes the positive side explicit: \emph{what an admitted
inside-view metaphor can deliver}. It shows that the description of
$T^4$ remains coherent when time is read not as a distinguished
continuous variable but as a further mode on the torus -- and that
the fractal structure that appears in FFGFT is, in this reading, not
a property of the underlying geometry but a phenomenon that arises in
the transition to decompactified time. Within this transition -- and
only there -- a particular inside-view metaphor becomes admissible,
which is named explicitly and framed carefully at the end.

% ==============================================================================
\section{Time as mode, not stage}
% ==============================================================================

On $T^4$ all four coordinates are equally periodic. None of them is
distinguished as ``time''; the division into ``three spatial dimensions
plus one temporal dimension'' is not laid down in the carrier itself,
but arises only from the choice of a descriptive perspective. In the
FFGFT reading, what we phenomenally experience as \emph{flowing time}
is a \emph{mode} in the sense of an oscillation or motion pattern on
the torus -- not a stage on which physics takes place.

This position is consistent with two points already established in
the corpus:

\begin{itemize}[leftmargin=2em,itemsep=4pt,topsep=4pt]
  \item In FFGFT, $\hbar$ is an SI conversion factor between energy
        and time measures, not an ontological primitive
        (Doc.~001a, Doc.~230 §\,Operational equivalence). The status
        of time itself is therefore not primitive either, but
        derived.
  \item The foundational relation $\tilde T \cdot m = 1$ is, on
        $T^4$, a strict reciprocity, not an open time coordinate.
        Reciprocity is a property of a closed structure, not of a
        flow.
\end{itemize}

\textbf{Consequence.} What we perceive as ``continuously flowing
time'' is the phenomenal resolution of a discrete torus mode into a
quasi-continuous variable. This resolution we call \emph{decompactification
of time}. It is not wrong -- it delivers exactly the familiar
predictions of standard physics -- but it is a \emph{descriptive
choice}, not a property of the carrier.

% ==============================================================================
\section{Borrowed time: $\tilde T$ and energy as reciprocal readings}
% ==============================================================================

Before treating the deformation of the carrier by mass, an
observation is in order that clarifies the status of $\tilde T$
itself -- and that leads into a very old didactic parallel.

\textbf{The historical parallel.} In medieval Europe, negative
numbers were rejected because no one could picture ``$-3$ apples''.
Three apples borrowed, on the other hand, were obvious even then.
Acceptance came not through a better picture but through a shift in
the descriptive frame: not \emph{inventory of apple objects} but
\emph{debt relation}. As soon as the frame shifted, the matter was
immediately intelligible (cf.\ Part~A §\,1).

\textbf{The transfer to $\tilde T$.} What in the apple case is the
shift from \emph{object inventory} to \emph{debt relation} is, in
the case of time, the shift from \emph{flowing time} to
\emph{borrowed energy}. On $T^4$, $\tilde T$ is not an open time
coordinate but a \emph{reciprocal reading of the energetic action}
on the carrier. The foundational relation
\begin{equation}
  \tilde T \cdot m = 1
\end{equation}
says exactly this: $\tilde T$ and $m$ are not two distinct
ontological quantities but \emph{two reciprocal readings of the same
fundamental action}. A short $\tilde T$ corresponds to a high
energy, a long $\tilde T$ to a low one -- carrying the same
information, in inverse scaling. This is not mathematically new (in
natural units, energy is simply reciprocal time), but in FFGFT it is
taken ontologically seriously: $\tilde T$ is \emph{borrowed energy},
not an independent ``something''.

\textbf{What follows from this.} All dynamical quantities on $T^4$
can, in this reading, be understood as different scaling axes of
\emph{a single energetic action}: time ($\tilde T$), mass-energy
($m$, $E$), momentum, frequency, wave number. They differ not in
their ontological category but in the choice of scaling axis used
to express them. What remains as non-energetic substrate is only
the \emph{topology} of $T^4$: the four cyclic identifications, the
periodicity, the closure. This topology is not energy -- it is what
energy can be carried on.

\textbf{The methodological point.} As with borrowed apples,
acceptance is a question of the descriptive frame, not of intuition.
No one can spatially picture ``borrowed time'' -- no more than
``$-3$ apples''. But the structural notion of a reciprocal debt
relation is, in both cases, clear and tractable.
$\tilde T \cdot m = 1$ is, on $T^4$, what the debt relation is in
the apple case: a structural relation that needs no visualisation
to function.

This reading prepares the following section: if $\tilde T$ and $m$
are reciprocal readings of the same energetic action, then the
``deformation of the carrier by mass'' is nothing other than a
local redistribution of this action over the geometry. Curvature
and energy are not two different things that interact; they are
two aspects of the same.

% ==============================================================================
\section{Mass as deformation of the carrier}
% ==============================================================================

In general relativity, mass is described as curvature of the
spacetime manifold. On $T^4$, this idea carries over structurally:
what GR reads as local curvature of an open manifold becomes, on the
closed carrier, a local \emph{deformation} of the torus geometry. The
global topology is preserved -- the four identifications continue to
close the carrier -- but the local metric is altered by mass.

Three remarks:

\begin{itemize}[leftmargin=2em,itemsep=4pt,topsep=4pt]
  \item In the limit of small deformation (local inertial frames),
        the deformed torus is indistinguishable from the Euclidean
        spacetime of GR -- GR is, in this reading, the linear
        approximation in which the global closure is invisible.
  \item GR singularities (black hole, Big Bang) appear structurally
        differently on $T^4$ because the carrier is closed through
        identifications (cf.\ Doc.~230 §\,Operational equivalence,
        the point on singularities).
  \item The deformation is not visualisable -- for exactly the
        reasons developed in Doc.~162. What is described here is the
        \emph{structural relationship} between mass and carrier, not
        a picture.
\end{itemize}

% ==============================================================================
\section{$T^4$ as a scale-structural, not scale-metric, statement}
% ==============================================================================

Up to this point $T^4$ has been treated as though it were \emph{one}
geometry of a fixed size. The consequences of the model are drawn,
however, on very different length scales: lepton masses
($\sim 10^{-18}\,\mathrm{m}$), Bell correlations (laboratory,
$\mathrm{m}$ to $\mathrm{km}$), material dispersion in
$\mathrm{LiNbO}_3$ ($\sim \mu\mathrm{m}$), Hubble scale and
$\Lambda$ ($\sim 10^{26}\,\mathrm{m}$). A critical reader is right to
ask: ``What scale does the torus have, then?''

\textbf{The answer.} $T^4$ is not an object with a fixed length
scale but a \emph{structural property} of the carrier geometry,
whose closure invariant $\xi$ is \emph{scale-invariant}. On every
length scale the same closure principle manifests with
scale-specific resolution. It does not follow that there are
``many tori'' -- it follows that the statement ``torus'' is, in
FFGFT, a \emph{structural} and not a \emph{metric} claim. What is
the same on every scale is the geometric closure; what differs from
scale to scale is the concrete resolution in which this closure
becomes visible.

\textbf{Empirical trace.} This scale-structural reading is not mere
stipulation. It is supported by the appearance of $\xi$ as
\emph{the same dimensionless constant} in entirely different
length contexts:

\begin{itemize}[leftmargin=2em,itemsep=2pt,topsep=2pt]
  \item lepton-mass cascade ($\xi$ in sub-femtometre physics),
  \item Higgs formula $\xi = \lambda_h^2 v^2 / (16\pi^3 m_h^2)$
        (electroweak sector),
  \item fine-structure constant $\alpha = e^2 \mu_0 c / (4\pi\hbar)$
        (atomic physics),
  \item $\mathrm{LiNbO}_3$ dispersion difference
        $\Delta n = \sqrt{3}\,\xi^{1/2}$ (material physics at
        \textmu m scale; Doc.~167, agreement to 16 significant
        digits),
  \item Koide relation $\xi_{\text{Koide}} \approx 1.372 \times
        10^{-4}$ (lepton ratios).
\end{itemize}

That the same number appears in five mutually independent physical
contexts at very different length scales is empirically non-trivial
and is exactly what a scale-invariant closure structure would lead
one to expect.

\textbf{Structural parallel to the renormalisation group.} Standard
physics already knows this pattern: renormalisation-group methods,
asymptotic freedom of QCD, effective field theories work with the
same idea -- that the same theoretical structure is effective at
different scales, with scale-dependent effective manifestations.
FFGFT is not outside this tradition here; it makes the pattern
explicit: $\xi$ is the scale-invariant quantity, the length scale
is the system-dependent resolution. This is also consistent with
Doc.~182, in which it is set out that FFGFT recognises \emph{no
universal maximum extent} -- each scale has a system-specific upper
bound.

\textbf{The methodological point.} This observation connects to the
reading of $\tilde T$ and energy as reciprocal readings developed
in §\,2: just as $\tilde T$ and $m$ are not two different
quantities but two scaling axes of the same action, the tori on
different length scales are not different structures but the same
structural property on different resolution axes. The universality
of $\xi$ is the operational manifestation of this identity.

% ==============================================================================
\section{Fractality as an inside-view phenomenon}
% ==============================================================================

FFGFT contains a fractal correction structure (factor $K_{\text{frak}}$,
appearing in $\pi$-gates and in the lepton-mass cascade). Until now
this fractality has been treated in the corpus as a geometric
property of $T^4$. The discussion of time-as-mode shows that a more
precise reading is possible:

\begin{quote}
As long as time is a discrete mode on the torus, there is no genuine
self-iteration. A fractal recursion requires the application of an
operation to its own result -- and that requires a direction in
which ``before'' and ``after'' can be distinguished. On a closed
torus mode this direction does not exist. It arises only when the
mode is decompactified into a directed, quasi-continuous time.
\end{quote}

In this reading, the fractal structure is not an ontological feature
of the underlying geometry but a \emph{phenomenon of the
decompactified time description}. It appears as soon as one enters
the mode ``directed time'' -- that is, as soon as one adopts an
inside view in which ``before'' and ``after'' can be distinguished.

\textbf{Structural parallel to the Born rule.} This position is not
isolated. It is structurally parallel to the Born-rule reading from
Doc.~230 (§\,Operational equivalence): there too, the apparent
randomness of quantum measurement arises only in the transition from
the deterministic $T^4$ dynamics to decompactified time. Both
phenomena -- apparent randomness and fractal inside-geometry -- are
consequences of the \emph{same decompactification}, not independent
findings.

This structural parallel is a non-trivial result: two phenomena that
stand unconnected in standard physics (quantum randomness and
fractal corrections) appear in FFGFT as two manifestations of a
single cause.

\textbf{The computational proof.} The reading above is not only an
ontological clarification; it is a \emph{computationally verifiable
result}. In natural units ($\hbar = c = 1$, dimensionless ratios),
the entire FFGFT description is \emph{correction-free} -- no fractal
corrections occur, neither $K_{\text{frak}}$ nor any other
adjustments. Only in the translation to SI units -- the moment in
which the energetic action is projected into dimensional quantities
with separated units (second, joule, metre) -- do the fractal
corrections appear. This is set out in detail in Doc.~261
(§\,Static geometry and SI projection): as long as one computes in
natural units and pure ratios, the description is static and
correction-free (Layer~1, compact $T^4$, frozen recursion). $c$,
$\hbar$, $\varepsilon_0$ become real only upon decompactification of
the intrinsic time winding $\tilde T$ into arrow time $t$
(Layer~2, embedding cost).

This shifts the epistemic status of the fractal correction: it is
\emph{not an additional term in the theory} that would have to be
added phenomenologically, but an \emph{artefact of the SI
projection} that arises precisely in the step in which the compact
$T^4$ description is translated into the decompactified SI
reading. Whoever computes in natural units stays in the Layer-1
description -- and sees no corrections. Whoever computes in SI
necessarily moves to the Layer-2 description -- and sees the
corrections, because they describe exactly what the projection
forces in terms of action.

This is the computational proof that fractality is a translation
phenomenon and not an ontological property of the carrier. The test
is simple and already carried out within the theory itself: the
same physical statement formulated in two unit systems. One shows
the corrections; the other does not. The distinguishing element is
not the physics but the \emph{descriptive choice}.

% ==============================================================================
\section{Inside-view metaphor: brain folds and light}
% ==============================================================================

\textbf{Preliminary note on admissibility.} The following metaphor is
introduced here \emph{explicitly}, because an openly named metaphor
remains controllable, while a hidden allusion invites suspicion of
esotericism. It is an \emph{inside-view metaphor}, not an ontological
claim about $T^4$. It does not say ``$T^4$ looks like a brain''; it
says something more specific and weaker.

\textbf{The metaphor.} Once time is decompactified, an observer
inside the geometry experiences not the simple structure of the
carrier but a fractally folded experiential geometry -- with folds,
fold-curvatures, and self-similarities on several scales. The
folded structure of a brain -- many-times wound, self-similar,
geometrically rich on a small volume -- is a familiar image for
exactly this kind of inside structure. It is not \emph{what} $T^4$
is, but it is an appropriate visualisation of what the
\emph{inside view of decompactified time} sees.

\textbf{The light example sharpens the point.} \emph{Viewed from
outside} -- in the undistorted description -- light moves in a
straight line. Within the deformed geometry, however, the path
actually traversed is different: it follows the folds and windings,
and for the observer situated inside, \emph{this} winding path is
the real one, not the idealised straight line of the outside
description. The standard reading -- straight light paths in
Euclidean space -- is, in this view, the \emph{outside description},
which holds in the limit of undistorted geometry; the inside view
is the description in which we actually live and think.

\textbf{The clean framing -- what the metaphor delivers and what
not.}

\begin{itemize}[leftmargin=2em,itemsep=2pt,topsep=2pt]
  \item \emph{What it delivers:} It illustrates that a fractally
        complex inside-experiential world is compatible with a
        simple outside carrier structure. It makes plausible why
        our phenomenal world \emph{appears} complex while its
        carrier is simple.
  \item \emph{What it is not:} It is not a claim that the universe
        is ``like a brain'', or that consciousness is ``fractal
        cosmic structure'', or anything similar. Such statements
        would overstep the frame from metaphor to ontology and
        would not be defensible within FFGFT.
  \item \emph{Where it applies:} Only in the transition ``discrete
        mode $\to$ decompactified time''. Outside this transition --
        in the pure $T^4$ description -- there is no fractal
        inside-geometry to describe. The metaphor applies to the
        inside view, not to $T^4$ itself.
\end{itemize}

Within this delimitation the metaphor is methodologically admissible
and didactically valuable. It closes a gap that Doc.~162 left open:
Doc.~162 showed what \emph{no} form of visualisation can deliver
(spatial intuition); this document shows what a \emph{deliberately
delimited} inside-view metaphor can deliver (rendering plausible the
relationship between a simple outside structure and a complex
inside world).

% ==============================================================================
\section{Consequences}
% ==============================================================================

From the four preceding sections, three statements follow that
individually already stand in the corpus, but are brought together
here for the first time as a common structural line:

\begin{enumerate}[leftmargin=2em,itemsep=4pt,topsep=4pt]
  \item \textbf{Time is not distinguished.} It is a mode among four
        equally periodic coordinates on $T^4$. Phenomenal
        time-experience is the inside view of decompactified time.

  \item \textbf{Mass is deformation of the carrier.} What GR
        describes as spacetime curvature is, on $T^4$, a local
        deformation of the closed geometry -- with standard GR as
        the linear, locally inertial approximation.

  \item \textbf{Fractality is an inside-view phenomenon.} The
        fractal structure that FFGFT sees appearing in
        $K_{\text{frak}}$, in the lepton-mass cascade, and in the
        $\pi$-gates is a phenomenon of decompactified time, not a
        property of the carrier itself. Parallel to the Born-rule
        reading from Doc.~230, two otherwise unconnected phenomena
        -- quantum randomness and fractal corrections -- thus
        appear as two manifestations of the same cause.
\end{enumerate}

\textbf{Methodological position.} This document tightens no
empirical prediction and changes no calculation. It is a
\emph{reading clarification}: the structural relationship between
the $T^4$ carrier, decompactified time, and phenomenal
experiential world is made explicit. The practical use is that the
metaphor of brain folds -- and similar inside-view images -- now
have a cleanly framed standing in the corpus and need no longer
operate as unexplained allusions.

% ==============================================================================
\section{Summary}
% ==============================================================================

On $T^4$, time is not a stage but a mode. Mass is not an object on
a spacetime but a deformation of the carrier. Fractality is not a
property of the carrier but a phenomenon arising in the transition
to decompactified time -- structurally parallel to the Born rule
and the apparent randomness of quantum measurement. Within this
transition the brain-folds metaphor is an admissible inside-view
visualisation -- not as a claim about $T^4$, but as an image for
how a simple outside structure becomes a fractally wound inside
world.

The outside description is simple. The inside world in which we
live is wound. The two descriptions are not in contradiction --
they are two sides of the same transition.

